%----------
%   IMPORTANTE
%----------

% Si nunca has utilizado LaTeX es conveniente que aprendas una serie de conceptos básicos antes de utilizar esta plantilla. Te aconsejamos que leas previamente algún tutorial (puedes encontar muchos en Internet).

% Esta plantilla está basada en las recomendaciones de la guía "Trabajo fin de Máster: Escribir el TFM", que encontrarás en http://uc3m.libguides.com/TFM/escribir
% contiene recomendaciones de la Biblioteca basadas principalmente en estilos APA e IEEE, pero debes seguir siempre las orientaciones de tu Tutor de TFM y la normativa de TFM para tu titulación.

% Encontrarás un ejemplo de TFM realizado con esta misma plantilla en la carpeta "_ejemplo_TFM_2019". Consúltalo porque contiene ejemplos útiles para incorporar tablas, figuras, listados de código, bibliografía, etc.



%----------
%	CONFIGURACIÓN DEL DOCUMENTO
%----------

% Definimos las características del documento y añadimos una serie de paquetes (\usepackage{package}) que agregan funcionalidades a LaTeX.

\documentclass[12pt]{report} %fuente a 12pt

%------------------------------------------------------------------------------------

\hyphenation{ciber-delincuentes pro-por-cionadas com-pro-meter ges-tio-nar la-bo-ra-torio} 



\usepackage{cite}


\usepackage{color}
\definecolor{gray97}{gray}{.97}
\definecolor{gray75}{gray}{.75}
\definecolor{gray45}{gray}{.45}

\usepackage{listings}
\lstset{ frame=Ltb,
framerule=0pt,
aboveskip=0.5cm,
framextopmargin=3pt,
framexbottommargin=3pt,
framexleftmargin=0.4cm,
framesep=0pt,
rulesep=.4pt,
backgroundcolor=\color{gray97},
rulesepcolor=\color{black},
%
stringstyle=\ttfamily,
showstringspaces = false,
basicstyle=\small\ttfamily,
commentstyle=\color{gray45},
keywordstyle=\bfseries,
%
numbers=left,
numbersep=15pt,
numberstyle=\tiny,
numberfirstline = false,
breaklines=true,
}

% minimizar fragmentado de listados
\lstnewenvironment{listing}[1][]
{\lstset{#1}\pagebreak[0]}{\pagebreak[0]}

\lstdefinestyle{consola}
{basicstyle=\scriptsize\bf\ttfamily,
backgroundcolor=\color{gray75},
}

\lstdefinestyle{C}
{language=C,
}

%------------------------------------------------------------------------------------



% MÁRGENES: 2,5 cm sup. e inf.; 3 cm izdo. y dcho.
\usepackage[
a4paper,
vmargin=2.5cm,
hmargin=3cm
]{geometry}

% INTERLINEADO: Estrecho (6 ptos./interlineado 1,15) o Moderado (6 ptos./interlineado 1,5)
\renewcommand{\baselinestretch}{1.15}
\parskip=6pt

% DEFINICIÓN DE COLORES para portada y listados de código
\usepackage[table]{xcolor}
\definecolor{azulUC3M}{RGB}{0,0,102}
\definecolor{gray97}{gray}{.97}
\definecolor{gray75}{gray}{.75}
\definecolor{gray45}{gray}{.45}

% Soporte para GENERAR PDF/A --es importante de cara a su inclusión en e-Archivo porque es el formato óptimo de preservación y a la generación de metadatos, tal y como se describe en http://uc3m.libguides.com/ld.php?content_id=31389625. En la carpeta incluímos el archivo plantilla_tfg_2017.xmpdata en el que puedes incluir los metadatos que se incorporarán al archivo PDF cuando lo compiles. Ese archivo debe llamarse igual que tu archivo .tex. Puedes ver un ejemplo en esta misma carpeta.
\usepackage[a-1b]{pdfx}

% ENLACES
\usepackage{hyperref}
\hypersetup{colorlinks=true,
	linkcolor=black, % enlaces a partes del documento (p.e. índice) en color negro
	urlcolor=blue} % enlaces a recursos fuera del documento en azul

% EXPRESIONES MATEMATICAS
\usepackage{amsmath,amssymb,amsfonts,amsthm}

\usepackage{txfonts} 
\usepackage[T1]{fontenc}
\usepackage[utf8]{inputenc}

\usepackage[spanish, es-tabla]{babel} 
\usepackage[babel, spanish=spanish]{csquotes}
\AtBeginEnvironment{quote}{\small}

% diseño de PIE DE PÁGINA
\usepackage{fancyhdr}
\pagestyle{fancy}
\fancyhf{}
\renewcommand{\headrulewidth}{0pt}
\rfoot{\thepage}
\fancypagestyle{plain}{\pagestyle{fancy}}

% DISEÑO DE LOS TÍTULOS de las partes del trabajo (capítulos y epígrafes o subcapítulos)
\usepackage{titlesec}
\usepackage{titletoc}

\titleformat{\chapter}[block]
{\large\bfseries} %\filcenter}
{\thechapter.}
{5pt}
{}
{}
\titlespacing{\chapter}{0pt}{0pt}{*3}
\titlecontents{chapter}
[0pt]                                               
{}
{\contentsmargin{0pt}\thecontentslabel.\enspace}
{\contentsmargin{0pt}}                        
{\titlerule*[.7pc]{.}\contentspage}                 

\titleformat{\section}
{\bfseries}
{\thesection.}
{5pt}
{}
\titlecontents{section}
[5pt]                                               
{}
{\contentsmargin{0pt}\thecontentslabel.\enspace}
{\contentsmargin{0pt}}
{\titlerule*[.7pc]{.}\contentspage}

\titleformat{\subsection}
{\normalsize\bfseries}
{\thesubsection.}
{5pt}
{}
\titlecontents{subsection}
[10pt]                                               
{}
{\contentsmargin{0pt}                          
	\thecontentslabel.\enspace}
{\contentsmargin{0pt}}                        
{\titlerule*[.7pc]{.}\contentspage}  


% DISEÑO DE TABLAS. Puedes elegir entre el estilo para ingeniería o para ciencias sociales y humanidades. Por defecto, está activado el estilo de ingeniería. Si deseas utilizar el otro, comenta las líneas del diseño de ingeniería y descomenta las del diseño de ciencias sociales y humanidades
\usepackage{multirow} % permite combinar celdas 
\usepackage{caption} % para personalizar el título de tablas y figuras
\usepackage{floatrow} % utilizamos este paquete y sus macros \ttabbox y \ffigbox para alinear los nombres de tablas y figuras de acuerdo con el estilo definido. Para su uso ver archivo de ejemplo 
\usepackage{array} % con este paquete podemos definir en la siguiente línea un nuevo tipo de columna para tablas: ancho personalizado y contenido centrado
\newcolumntype{P}[1]{>{\centering\arraybackslash}p{#1}}
\DeclareCaptionFormat{upper}{#1#2\uppercase{#3}\par}

% Diseño de tabla para ingeniería
\captionsetup[table]{
	format=upper,
	name=TABLA,
	justification=centering,
	labelsep=period,
	width=.75\linewidth,
	labelfont=small,
	font=small,
}

%Diseño de tabla para ciencias sociales y humanidades
%\captionsetup[table]{
%	justification=raggedright,
%	labelsep=period,
%	labelfont=small,
%	singlelinecheck=false,
%	font={small,bf}
%}


% DISEÑO DE FIGURAS. Puedes elegir entre el estilo para ingeniería o para ciencias sociales y humanidades. Por defecto, está activado el estilo de ingeniería. Si deseas utilizar el otro, comenta las líneas del diseño de ingeniería y descomenta las del diseño de ciencias sociales y humanidades
\usepackage{graphicx}
\graphicspath{{imagenes/}} %ruta a la carpeta de imágenes

% Diseño de figuras para ingeniería
\captionsetup[figure]{
	format=hang,
	name=Fig.,
	singlelinecheck=off,
	labelsep=period,
	labelfont=small,
	font=small		
}

% Diseño de figuras para ciencias sociales y humanidades
%\captionsetup[figure]{
%	format=hang,
%	name=Figura,
%	singlelinecheck=off,
%	labelsep=period,
%	labelfont=small,
%	font=small		
%}


% NOTAS A PIE DE PÁGINA
\usepackage{chngcntr} %para numeración contínua de las notas al pie
\counterwithout{footnote}{chapter}

% LISTADOS DE CÓDIGO
% soporte y estilo para listados de código. Más información en https://es.wikibooks.org/wiki/Manual_de_LaTeX/Listados_de_código/Listados_con_listings
\usepackage{listings}

% definimos un estilo de listings
\lstdefinestyle{estilo}{ frame=Ltb,
	framerule=0pt,
	aboveskip=0.5cm,
	framextopmargin=3pt,
	framexbottommargin=3pt,
	framexleftmargin=0.4cm,
	framesep=0pt,
	rulesep=.4pt,
	backgroundcolor=\color{gray97},
	rulesepcolor=\color{black},
	%
	basicstyle=\ttfamily\footnotesize,
	keywordstyle=\bfseries,
	stringstyle=\ttfamily,
	showstringspaces = false,
	commentstyle=\color{gray45},     
	%
	numbers=left,
	numbersep=15pt,
	numberstyle=\tiny,
	numberfirstline = false,
	breaklines=true,
	xleftmargin=\parindent
}

\captionsetup[lstlisting]{font=small, labelsep=period}
% fijamos el estilo a utilizar 
\lstset{style=estilo}
\renewcommand{\lstlistingname}{\uppercase{Código}}


%BIBLIOGRAFÍA - PUEDES ELEGIR ENTRE ESTILO IEEE O APA. POR DEFECTO ESTÁ CONFIGURADO IEEE. SI DESEAS USAR APA, COMENTA LAS LÍNEA DE IEEE Y DESCOMENTA LAS DE APA. Si haces cambios en la configuración de la bibliografía y no obtienes los resultados esperados, es recomendable limpiar los archivos auxiliares y volver a compilar en este orden: COMPILAR-BIBLIOGRAFIA-COMPILAR

% Tienes más información sobre cómo generar bibliografía y CONFIGURAR TU EDITOR DE TEXTO para compilar con biber en http://tex.stackexchange.com/questions/154751/biblatex-with-biber-configuring-my-editor-to-avoid-undefined-citations , https://www.overleaf.com/learn/latex/Bibliography_management_in_LaTeX y en http://www.ctan.org/tex-archive/macros/latex/exptl/biblatex-contrib
% También te recomendamos consultar la guía temática de la Biblioteca sobre citas bibliográficas: http://uc3m.libguides.com/guias_tematicas/citas_bibliograficas/inicio

% CONFIGURACIÓN PARA LA BIBLIOGRAFÍA IEEE
% \usepackage[backend=biber, style=ieee, isbn=false,sortcites, maxbibnames=5, minbibnames=1]{biblatex} % Configuración para el estilo de citas de IEEE, recomendado para el área de ingeniería. "maxbibnames" indica que a partir de 5 autores trunque la lista en el primero (minbibnames) y añada "et al." tal y como se utiliza en el estilo IEEE.

%CONFIGURACIÓN PARA LA BIBLIOGRAFÍA APA
%\usepackage[style=apa, backend=biber, natbib=true, hyperref=true, uniquelist=false, sortcites]{biblatex}
%\DeclareLanguageMapping{spanish}{spanish-apa}

% Añadimos las siguientes indicaciones para mejorar la adaptación de los estilos en español
% \DefineBibliographyStrings{spanish}{%
% 	andothers = {et\addabbrvspace al\adddot}
% }
% \DefineBibliographyStrings{spanish}{
% 	url = {\adddot\space[En línea]\adddot\space Disponible en:}
% }
% \DefineBibliographyStrings{spanish}{
% 	urlseen = {Acceso:}
% }
% \DefineBibliographyStrings{spanish}{
% 	pages = {pp\adddot},
% 	page = {p.\adddot}
% }

% \addbibresource{Bibliografia/Bibliografia.bib} % llama al archivo bibliografia.bib en el que debería estar la bibliografía utilizada


%-------------
%	DOCUMENTO
%-------------

\begin{document}
\pagenumbering{roman} % Se utilizan cifras romanas en la numeración de las páginas previas al cuerpo del trabajo
	

%----------
%	PORTADA
%----------	

\thispagestyle{empty}


%----------
%	PORTADA
%----------	
\begin{titlepage}
	\begin{sffamily}
	\color{azulUC3M}
	\begin{center}
		\begin{figure}[H] %incluimos el logotipo de la Universidad
			\makebox[\textwidth][c]{\includegraphics[width=16cm]{Portada_Logo.png}}
		\end{figure}
		\vspace{1cm}
		\begin{Large}
			Máster Universitario en Ciberseguridad\\			
			2018-2019\\
			\vspace{1cm}		
			\textsl{Trabajo Fin de Máster}
			\bigskip
			
		\end{Large}
		 	% {\Huge ``CIBEREJERCICIOS PARA EVALUAR ACTIVE DIRECTORY EN SUS DISTINTAS VERSIONES. ''}\\
		 	{\Huge ``Ciberejercicios para evaluar \\
		 	Active Directory en sus distinas versiones. ''}\\
		 	\vspace*{0.5cm}
	 		\rule{10.5cm}{0.1mm}\\
			\vspace*{0.9cm}
			{\LARGE Borja Lorenzo Fernádez}\\ 
			\vspace*{1cm}
		\begin{Large}
			Tutor\\
			Andrés Marín López\\
		\end{Large}
	\end{center}
	
	\vfill
	\color{black}
	\begin{footnotesize}
	\noindent\fbox{
	\begin{minipage}{\textwidth}
		\textbf{DETECCIÓN DEL PLAGIO}\\
		La Universidad utiliza el programa \textbf{Turnitin Feedback Studio} para comparar la originalidad del trabajo entregado por cada estudiante con millones de recursos electrónicos y detecta aquellas partes del texto copiadas y pegadas. Copiar o plagiar en un TFM es considerado una \textbf{\underline{Falta Grave}}, y puede conllevar la expulsión definitiva de la Universidad.
	\end{minipage}	
	}
	% \vspace*{.5cm}\\	
	% \noindent\includegraphics[width=4.2cm]{imagenes/creativecommons.png}\\
	% \emph{[Incluir en el caso del interés en su publicación en el archivo abierto]}\\
	% Esta obra se encuentra sujeta a la licencia Creative Commons \textbf{Reconocimiento - No Comercial - Sin Obra Derivada}

	\end{footnotesize}
	\end{sffamily}
\end{titlepage}

\newpage %página en blanco o de cortesía
\thispagestyle{empty}
\mbox{}

%----------
%	Agracedimientos
%----------	
\chapter*{Agradecimientos}
%----------
%	Agracedimientos
%----------	

\setcounter{page}{3}
	
	% ESCRIBIR LA DEDICATORIA AQUÍ	


En primer lugar, me gustaría dedicar este trabajo a mis abuelos ya que sin ellos no habría podido llegar a donde he llegado. \\

Agradecer especialmente a mi familia, padres y hermanos por ayudarme y aguantarme incluso en los peores momentos. \\

Quiero agradecer a María, por la dedicación, el esfuerzo y la paciencia que ha tenido conmigo y por habeme enseñado tanto y haberme facilitado las cosas para llegar hasta aquí. \\

Por último, agradecer a todos los que me han ayudado con el desarrollo de este proyecto desinteresadamente, especialemente a Atl4s, Arcocapaz, Cynops y Mamatb.
\vfill
\newpage %página en blanco o de cortesía
\thispagestyle{empty}
\mbox{}


%----------
%	RESUMEN Y PALABRAS CLAVE
%----------	
\chapter*{Resumen}
%----------
%	RESUMEN Y PALABRAS CLAVE
%----------	

\setcounter{page}{5}
	
Active Directory (AD) es una base de datos distribuida que permite a los administradores gestionar de manera eficiente la información de la empresa y los límites de la misma. En la actualidad, Active Directory es considerado una de las partes fundamentales para el correcto funcionamiento de cualquier empresa. La información gestionada por Active Directory, engloba usuarios, clientes, proveedores, dispositivos (ordenadores, servidores, impresoras...), servicios y aplicaciones así como la interacción entre ellos. Por lo tanto, supone uno de los principales objetivos para atacantes y ciberdelincuentes así como una gran problema de seguridad para una organización si el sistema de gestión principal se ve comprometido. En los últimos años se ha producido un gran aumento de ataques contra Active Directory cuyo objetivo es hacerse con el control de la empresa y comprometer la seguridad de la información, esto es debido a que se han perfeccionado los principales ataques así como la implementación de nuevas herramientas que facilitan todo tipo de ataques. \\

Con el fin de contribuir al desarrollo y mejora de la seguridad de Active Directory y por ende, de las empresas y particulares que hacen uso de ello, el trabajo realizado se centra en el estudio de los principales componentes que engloban la seguridad y gestión del sistema de directorio Active Directory a través de la creación de un laboratorio, de manera local, y el análisis y revisión de los principales ataques y amenazas usadas en la actualidad para vulnerar dicho sistema de gestión. Para ello, se implementa una topología que simula una empresa ficticia e implementa la última versión de Windows Server 2019.\\ 

\textit{\textbf{Palabras Clave:}}

Microsoft Active Directory, Domain Controller, Kerberos, Ciberseguridad, Windows Server, Pentesting, Red Team
	
		
\vfill
\newpage %página en blanco o de cortesía
\thispagestyle{empty}
\mbox{}






	

%-------------------------------------------------------------------------------
%	ÍNDICES
%-------------------------------------------------------------------------------

%--
%Índice general
%-
\tableofcontents
\thispagestyle{fancy}

\newpage %página en blanco o de cortesía
\thispagestyle{empty}
\mbox{}

%--
%Índice de figuras. Si no se incluyen, comenta las líneas siguientes
%-
\listoffigures
\thispagestyle{fancy}

\newpage %página en blanco o de cortesía
\thispagestyle{empty}
\mbox{}

%--
%Índice de tablas. Si no se incluyen, comenta las líneas siguientes
%-
\listoftables
\thispagestyle{fancy}

\newpage %página en blanco o de cortesía
\thispagestyle{empty}
\mbox{}


%-------------------------------------------------------------------------------
%-------------------------------------------------------------------------------

%----------
%	TRABAJO
%----------	
\clearpage
\pagenumbering{arabic} % numeración con múmeros arábigos para el resto de la publicación	


\chapter{Introducción}
En los últimos años, empresas y organizaciones se han visto en la necesidad de gestionar de una manera eficiente y centralizada la información y recursos en red que disponen, activo fundamental para el correcto funcionamiento del negocio. El aumento masivo de dicha información además de la necesidad de crear, distribuir y manipular tal cantidad de datos, ya sea a través de servicios de bases de datos como puede ser servicios MySQL, la obtención de servidores y dispositivos para su almacenamiento, la creación de  aplicaciones web y servicios que permitan su distribución o la gestión de los usuarios que puedan manipularla o consultarla supone una gran cantidad de agentes que intervienen en el funcionamiento que es necesario controlar y regular. Microsoft Active Directory supone una solución a esa problemática a través de un servicio de directorio como base de datos distribuida que permite la gestión, administración y localización de todos los recursos en red ~\cite{Capitulo1:Microsoft}.\\

Dentro del panorama actual, los servicios de Microsoft Active Directory se ha convertido en uno de los pilares que sostienen la organización de los recursos en red de la mayoría de las empresas vigentes así como uno de los principales objetivos para atacantes  debido a dicha importancia. Esto se puede comprobar en el principal interés que tienen los principales grupos de ciberdelincuentes como APT28, Cobalt Strike...  por esta infraestructura o los ataques de ransomware WannaCry, NotPetya, MBR-ONI, etc, que ponen a Active Directory en el punto de mira y centro de sus ataques \cite{Capitulo1:Ransomware}. En los últimos años, se ha visto un aumento considerable de vulnerabilidades críticas que afectan a la seguridad y que detectan y hacen considerablemente más sencillo su explotación. \\

Por todo ello y con el fin de abordar esta problemática, el trabajo realizado se ha centrado en la revisión, análisis y prueba en profundidad de las principales amenazas que suponen un problema de seguridad para Active Directory en sus diferentes versiones, técnicas como Pass-The-Hass, NTLM Relay, Kerberoast, etc que serán detalladas en los capítulos posteriores. Además, también se proporciona las directrices para la creación de un laboratorio local que permite la prueba de los ataques detallados así como la posibilidad de probar nuevas técnicas y ataques sin poner en riesgo la seguriidad de ningún entorno real u organización.\\


\section{Estado del Arte}

Para el desarrollo del trabajo, se ha considerado las siguientes versiones proporcionadas por Microsoft para la instalación de los dominios que van a formar partes del Active Directory y van a ejecutar los Domain Controllers:\\

\begin{itemize}
\item \textbf{Windows Server 2019}
\item \textbf{Windows Server 2016}
\item \textbf{Windows Server 2012 R2}
\end{itemize}

Por lo tanto, se ha dejado la versión más antigua Windows Server 2008 como objeto de estudio o posible implantación en trabajo futuro al ser la versión más obsoleta. Aunque esto no implica que no haya empresas que aún siguen usando dicha versión. \\

En cuanto a los ataques a analizar de manera detallada se han considerado los siguientes: 

\begin{itemize}
\item \textbf{Pass-The-Hash}
\item \textbf{NTLM Relay}
\item \textbf{Overpass-The-Hash}
\item \textbf{Pass-The-Ticket}
\item \textbf{Golden/Silver Ticket}
\item \textbf{Kerberoast}
\end{itemize}

Como se puede observar, la mayoría de los ataques no son específicos de Active Directory si no que atacan a los protocolos NTLM y Kerberos, por lo que, también se van a detallar en profundidad estos protocolos en los capitulos siguientes.


\section{Objetivos}

El objetivo principal de este trabajo es la creación de un laboratorio que provea la infraestructura necesaria para la replicación de las principales técnicas de ataque a Active Directory así como la revisión de las mismas sobre las diferentes versiones proporcionadas por Microsoft. La creación del laboratorio posibilita tanto a {\it Pentesters} o expertos en seguridad ofensiva realizar ejercicios de {\it Read Team} en entornos controlados o réplicas de un entorno real como a administradores de sistemas para probar nuevas configuraciones y reglas para equipos de {\it Blue Team}.\\

Además, este trabajo tiene como objetivo la adquisición de conocimiento sobre Active Directory como punto de partida para futuras investigaciones. Como se ha visto en la introducción Active Directory es una parte fundamental de una empresa y es uno de los principales objetivos de atacantes, conocer los principales ataques y cómo está organizado es de gran importancia hoy en día, permitiendo aasí una correcta implementación que minimice los riesgos a los que está sometido.\\


\section{Organización del Proyecto}

El presente documento se divide en 7 capítulos, en los cuales en primera instancia se detallan los aspectos a tener en cuenta relacionados con Active Directoy, se detalla el laboratorio implementado, las pruebas que se van a realizar, la experimentación realizada así como los resultados obtenidos durante el trascurso:\\

En el Capítulo 2 se detalla los aspectos relacionados con Active Directory, en primer lugar se define la autenticación y autorización en sistemas Windows tanto localmente como en dominio, se analizan los protocolos NT Lan Manager y Kerberos y la terminología relacionada con Active Directory.\\

En el Capitulo 3 se desarrolla la topología que se ha elegido para la creación del laboratorio de pruebas con las diferentes versiones de Windows y su implementación.\\

En el Capítulo 4 se detallan los ataques elegidos para realizar la experimentación.\\

En el Capítulo 5, una vez detalladas tanto el laboratorio como las ataaques principales objeto de estudio, se muestran las diferentes pruebas realizas en las versiones de Windows especificadas en el estado del arte.\\

En el Capítulo 6 se presentan y se discuten los resultados obtenidos.\\

En el Capítulo 7, para finalizar el proyecto, se lleva a cabo una reflexión sobre el esfuerzo realizado y sus diferentes líneas de trabajo futuro. \\


\chapter{Aspectos Clave}
\section{Autenticación}

\section{NT Lan Managey (NTLM) y Kerberos}

\section{Active Directory}






%----------
%	BIBLIOGRAFÍA
%----------	

%\nocite{*} % Si quieres que aparezcan en la bibliografía todos los documentos que la componen (también los que no estén citados en el texto) descomenta está lína

\clearpage
\addcontentsline{toc}{chapter}{Bibliografía}
%\setquotestyle[english]{british} % Cambiamos el tipo de cita porque en el estilo IEEE se usan las comillas inglesas.
%\printbibliography

\bibliography{Bibliografia/Bibliografia}{}
\bibliographystyle{ieeetr}


%----------
%	ANEXOS
%----------	

% Si tu trabajo incluye anexos, puedes descomentar las siguientes líneas
%\chapter* {Anexo x}
%\pagenumbering{gobble} % Las páginas de los anexos no se numeran



\end{document}
